El desarrollo de la plataforma web centralizada para la Universidad del Valle de Guatemala (UVG) logró demostrar
la viabilidad de integrar con éxito flujos administrativos manuales y comunicación fragmentada en un sistema
digitalizado. A través de la implementación de una ecosistema digital unificado, se comprobó que es posible
gestionar el ciclo de vida de los proyectos de extensión, contemplando desde la publicación, hasta la verificación
de horas de trabajo reportada por estudiantes, de manera trazable y segura; cumpliendo a cabalidad la administración,
publicación, participación y seguimiento de los proyectos de extensión de la UVG, que permita a estudiantes,
directores y organizaciones externas participar oportunamente en el ciclo de vida de los proyectos de extensión.

Se logró implementar una API modular y escalable, mediante la selección de tecnologías modernas no bloqueantes,
una arquitectura por capas y un diseño de base de datos relacional, permitió demostrar un nivel óptimo de 
escalabilidad, rendimiento y modularidad. Las pruebas de estrés confirmaron que la API fue capaz de soportar
cargas de hasta 250 solicitudes concurrentes manteniendo tiempos de respuesta promedio inferiores a 300 ms,
lo que indica que puede soportar un número significativo de usuarios simultáneos sin degradar la experiencia
del usuario. Asimismo, la definición de reglas de negocio claras y la implementación de un sistema de
autenticación y autorización mediante JWT, garantizó la protección, control de acceso y seguridad de los
datos, asegurando que solo usuarios autorizados puedan realizar acciones específicas dentro de la plataforma.

Se materializó una interfaz de usuario intuitiva y responsiva a través de una \textit{Single Page Application}
(SPA), basada en los prototipos previamente diseñados y aprobados, para que directores, estudiantes y
organizaciones externas realicen sus gestiones sobre los proyectos de extensión. 
La integración de validaciones síncronas y manejo de caché de los datos obtenidos de la API, permitió optimizar
la experiencia del usuario.

Se verificó que la API brinda la funcionalidad necesaria para gestionar el ciclo de vida de los proyectos de
extensión, así como la capacidad de manejar otras funcionalidades necesarias para la administración, publicación,
participación y seguimiento de los proyectos de extensión de la UVG.
Además, las pruebas de usabilidad realizadas con usuarios finales confirmaron empíricamente que la plataforma podría
sustituir los flujos administrativos manuales y comunicación fragmentada, logrando una experiencia de usuario
satisfactoria, con una tasa de finalización del 100\% en las tareas asignadas durante la prueba, la mayoría
completadas en tiempos inferiores a dos minutos. La evaluación cuantitativa mediante la Escala de Usabilidad
del Sistema (SUS) arrojó un puntaje global de 85, lo que indica que la plataforma es altamente usable y bien
recibida por los usuarios finales.

